\subsection{取值于测度空间的函数}

设 X 是一个局部紧空间，\(\mathcal{M}_{+}(X)\) 是 X 上正测度所成的凸锥。在本章其余部分，\(\mathcal{M}_{+}(X)\) 总赋有由 \(\mathcal{M}(X)\) 上的淡拓扑诱导的拓扑（第三章，§1 第 9 号）；于是，说映射 \(\Lambda: t \mapsto \lambda_{t}\) （局部紧空间 T 到 \(\mathcal{M}_{+}(X)\) 中的）连续，意思是对每个函数 \(f \in \mathcal{K}(X)\) ，数值函数 \(t \mapsto \lambda_{t}(f)\) 连续。这时我们也称 \(\Lambda\) 在 T 上淡连续。说映射 \(\Lambda: t \mapsto \lambda_{t}\) 是 \(\mu\) -可测的，意思是 T 的紧子集 K （使 \(\Lambda\) 在 K 上的限制淡连续者）所成之集是 \(\mu\) -稠密的（第四章，§5 第 10 号命题 15）。这时我们将称 \(\Lambda\) 是淡 \(\mu\) -可测的。

设 \(\Lambda : t \mapsto \lambda_t\) 是 T 到 \(\mathcal{M}_+(\mathrm{X})\) 中的一个映射；我们就称 \(\Lambda\) 对测度 \(\mu\) 是标量本质可积的，如果对每个函数 \(f \in \mathcal{K}(\mathrm{X})\) ，函数 \(t \mapsto \lambda_t(f)\) 是本质 \(\mu\) -可积的。若令 \(\nu(f) = \int \lambda_t(f) d\mu(t)\) ，则显然 \(\nu\) 是 \(\mathcal{K}(\mathrm{X})\) 上的一个正线性形式，从而是 X 上的一个测度（第三章，§1 第 5 号定理 1）。我们将称 \(\nu\) 为函数 \(\Lambda\) （取值于 \(\mathcal{M}_+(\mathrm{X})\) ）的积分，并写作 \(\nu = \int \lambda_t d\mu(t)\) 。

前面的定义是弱积分概念的一个特例，弱积分将在第六章中以一般方式处理。

若 \(f\) 表示 \(\mathcal{K}(\mathrm{X})\) 的一个元素，则积分 \(\int \lambda_t(f) d\mu(t)\) 滥用记号也可记作 \(\int d\mu(t) \int f(x) d\lambda_t(x)\) ；于是积分 \(\nu = \int \lambda_t d\mu(t)\) 的定义可写成

\[
\int f (x) d \nu (x) = \int d \mu (t) \int f (x) d \lambda_ {t} (x).\tag{1}
\]

在后面，对上积分、本质上积分以及取值于 Banach 空间的函数的积分，我们将作类似的记号滥用。

例. --- 1) 设 T 是一个离散空间，且 \(\mu\) 是在 T 的每个点放置质量 \(+1\) 所定义的 T 上的测度（第三章，§1 第 3 号）。设 \(h\) 是定义在 T 上的一个 \(\geqslant 0\) 的函数；由于函数 \(h\) 在 T 上是下半连续的（甚至是连续的），\(\mu^{*}(h) = \mu^{\bullet}(h) = \sum_{t\in \mathrm{T}}h(t)\) （第四章，§1 第 1 号例）。

对测度 \(\mu\) 而言，可积函数与本质可积函数的概念因此是同一的。在此情形下，说映射 \(t \mapsto \lambda_t\) （T 到 \(\mathcal{M}_+(X)\) 中的）是标量本质 \(\mu\) -可积的，等于说族 \((\lambda_t)_{t \in \mathrm{T}}\) 是可和的（§2 第 1 号），且这时有 \(\int \lambda_t d\mu(t) = \sum_{t \in \mathrm{T}} \lambda_t\) 。注意映射 \(t \mapsto \lambda_t\) 是淡连续的。

\begin{enumerate}
\def\labelenumi{\arabic{enumi})}
\setcounter{enumi}{1}
\tightlist
\item
  映射 \(t \mapsto \varepsilon_t\) （\(T\) 到 \(\mathcal{M}_+(T)\) 中的）是淡连续的、标量本质 \(\mu\) -可积的（\(\mu\) 为 \(T\) 上任一正测度），且有 \(\int \varepsilon_t d\mu(t) = \mu\) 。
\end{enumerate}

命题 1. --- 设 \(\mu\) 是 T 上正测度的一个递增有向族 \((\mu_i)_{i \in I}\) 的上确界；为使 \(\Lambda : t \mapsto \lambda_t\) 是标量本质 \(\mu\) -可积的，必须且只须 \(\Lambda\) 是标量本质 \(\mu_i\) -可积的（对一切 \(i \in I\) ），且族 \(\left( \int \lambda_t d\mu_i(t) \right)_{i \in I}\) 在 \(\mathcal{M}(X)\) 中有上界。这时有，

\[
\int \lambda_ {t} d \mu (t) = \sup _ {i \in \mathrm{I}} \int \lambda_ {t} d \mu_ {i} (t).\tag{2}
\]

因为，验证 \(\Lambda\) 对 T 上一个正测度 \(\eta\) 是标量本质可积的，归结为验证 \(t \mapsto \lambda_{t}(g)\) 是 \(\eta\) -可测的且关于 \(\eta\) 有有限本质上积分（对每个函数 \(g \in \mathcal{K}_{+}(X)\) ）。因而命题由 §1 第 4 号命题 11 及其推论 2 立即得出。

推论. --- 设 \(\mu\) 是正测度的一个可和族 \((\mu_{\alpha})_{\alpha \in \mathbf{A}}\) （\(\mathbf{T}\) 上的）之和；为使 \(\Lambda : t \mapsto \lambda_t\) 是标量本质 \(\mu\) -可积的，必须且只须 \(\Lambda\) 是标量本质 \(\mu_{\alpha}\) -可积的（对每个 \(\alpha \in \mathbf{A}\) ），且测度族 \(\int \lambda_t d\mu_{\alpha}(t)\) 可和。这时有

\[
\int \lambda_ {t} d \mu (t) = \sum_ {\alpha \in A} \int \lambda_ {t} d \mu_ {\alpha} (t).\tag{3}
\]

由此立即得出，每个标量本质 \(\mu\) -可积的映射也都是标量本质 \(\mu'\) -可积的（对每个测度 \(\mu' \leqslant \mu\) ）。

在本节中，我们将限于研究 T 到 \(\mathcal{M}_{+}(X)\) 中的具有下述定义所考虑的性质的标量本质可积映射：

定义 1. --- 设 X 是一个局部紧空间，\(\Lambda : t \mapsto \lambda_t\) 是一个标量本质 \(\mu\) -可积映射（T 到 \(\mathcal{M}_{+}(X)\) 中的），\(\nu\) 是 \(\Lambda\) 的积分。

我们称 \(\Lambda\) 是 \(\mu\) -预适宜的，如果对每个定义在 X 上的下半连续函数 \(f \geqslant 0\) ，函数 \(t \mapsto \int^{\bullet} f d\lambda_t\) 在 T 上 \(\mu\) -可测，且

\[
\int^ {\bullet} f (x) d \nu (x) = \int^ {\bullet} d \mu (t) \int^ {\bullet} f (x) d \lambda_ {t} (x).\tag{4}
\]

我们称 \(\Lambda\) 是 \(\mu\) -适宜的 (*)，如果 \(\Lambda\) 是 \(\mu'\) -预适宜的（对每个正测度 \(\mu' \leqslant \mu\) ）。

可以证明，若 \(\Lambda\) 是 \(\mu\) -预适宜的且测度 \(\nu = \int \lambda_{t} d\mu(t)\) 是适度的——特别地，若 X 在无穷远处可数——则 \(\Lambda\) 是 \(\mu\) -适宜的（习题 7）；然而，尚不知道这两个概念一般是否等价。在下面第 2 与第 3 号的命题叙述中，冠以 a) 或 b) 的断言立即推广到预适宜映射，而冠以 c) 的断言只对适宜映射成立。

下述命题常常使人能验证一个给定的映射是 \(\mu\) -适宜的。

命题 2. --- 设 \(\Lambda : t \mapsto \lambda_t\) 是一个标量本质 \(\mu\) -可积映射（T 到 \(\mathcal{M}_{+}(X)\) 中的），并设 \(\nu = \int \lambda_t d\mu(t)\) 。

\begin{enumerate}
\def\labelenumi{\alph{enumi})}
\tightlist
\item
  若 \(\Lambda\) 是淡连续的，则映射 \(t \mapsto \lambda_t^\bullet(f)\) 对每个定义在 X 上的下半连续函数 \(f \geqslant 0\) 是下半连续的，\(\Lambda\) 是 \(\mu\) -适宜的，且我们有关系式
\end{enumerate}

\[
\int^ {*} f (x) d \nu (x) = \int^ {*} d \mu (t) \int^ {*} f (x) d \lambda_ {t} (x).\tag{5}
\]

\begin{enumerate}
\def\labelenumi{\alph{enumi})}
\setcounter{enumi}{1}
\item
  若 \(\Lambda\) 是淡 \(\mu\) -可测的，则 \(\Lambda\) 是 \(\mu\) -适宜的。
\item
  若 X 的拓扑有可数基，则 \(\Lambda\) 是淡 \(\mu\) -可测的（从而也是 \(\mu\) -适宜的）。
\end{enumerate}

设 \(f\) 是一个 \(\geqslant 0\) 的下半连续函数（定义在 \(X\) 上）。设 \(F\) 是一个对关系 \(\leqslant\) 有向的集，由函数 \(g \in \mathcal{H}(X)\) （满足 \(0 \leqslant g \leqslant f\) ）组成。对 \(g \in F\) ，用 \(h_g\) 表示定义在 \(T\) 上由 \(h_g(t) = \lambda_t(g)\) 给出的函数。类似地，令

\[
h _ {f} (t) = \lambda_ {t} ^ {*} (f) = \lambda_ {t} ^ {\bullet} (f) = \sup _ {g \in \mathrm{F}} h _ {g} (t)
\]

（§1 第 1 号命题 4）。我们作下述比 a) 弱的假设：只设 \(\Lambda\) 在 S 上的限制是淡连续的，其中 S 是 T 的一个包含 \(\mu\) 的支集的闭子集。对 \(g \in \mathbf{F}\) ，用 \(\overline{h}_g\) 表示这样一个数值函数：它在 S 上与 \(h_g\) 相同，且取值 \(+\infty\) 于 \(\mathbb{C}S\) 上。令 \(\overline{h}_f = \sup_{g \in \mathbb{F}} \overline{h}_g\) ；则在 S 上 \(\overline{h}_f = h_f\) 。对每个 \(g \in \mathbf{F}\) ，函数 \(\overline{h}_g\) 是下半连续的；因而 \(\overline{h}_f\) 是下半连续的，且由于族 \((\overline{h}_g)_{g \in \mathbb{F}}\) 是有向的，

\[
\mu^ {*} (\overline {{h}} _ {f}) = \sup _ {g \in \mathrm{F}} \mu^ {*} (\overline {{h}} _ {g}) = \sup _ {g \in \mathrm{F}} \mu^ {*} (h _ {g}) = \sup _ {g \in \mathrm{F}} \nu (g) = \nu^ {*} (f)
\]

（第四章，§1 第 1 号定理 1 及 §2 第 3 号命题 6）。由于在 S 上 \(h_f = \overline{h}_f\) ，从而几乎处处如此，这也可写成 \(\mu^*(h_f) = \nu^*(f)\) ，这是一个与 (5) 恒等的等式。类似地，由于 \(f\) 与 \(\overline{h}_f\) 都是下半连续的，前面的关系式给出等式 \(\mu^\bullet(\overline{h}_f) = \nu^\bullet(f)\) （§1 第 1 号命题 4）；由于在 S 上 \(\overline{h}_f = h_f\) ，由此得出 \(\mu^\bullet(h_f) = \nu^\bullet(f)\) （§1 第 1 号命题 1），这是一个与 (4) 恒等的等式。因此映射 \(\Lambda\) 是 \(\mu\) -预适宜的；但在此论证中处处可把 \(\mu\) 换成 \(\mu' \leqslant \mu\) ，把 \(\nu\) 换成 \(\nu' = \int \lambda_t d\mu'(t)\) ，因为 \(\Lambda\) 也是标量本质 \(\mu'\) -可积的，且 S 包含 \(\mu'\) 的支集。由此得出 \(\Lambda\) 是 \(\mu\) -适宜的。

设 \(\Lambda\) 是淡连续的；可取 \(S = T\) ；则 \(h_f = \overline{h}_f\) 是下半连续的，这就完成了命题 a) 部分的证明。

设 \(\Lambda\) 是淡 \(\mu\) -可测的，我们来证明 b)。集 \(\mathfrak{K}\) （由 T 的紧子集 K （使 \(\Lambda\) 在 K 上的限制连续者）组成）是 \(\mu\) -稠密的（第四章，§5 第 10 号命题 15），故存在 T 上测度的一个可和族 \((\mu_{\alpha})_{\alpha \in \mathrm{A}}\) ，使 \(\mu = \sum_{\alpha \in \mathrm{A}} \mu_{\alpha}\) 且诸测度 \(\mu_{\alpha}\) 中每一个的支集都属于 \(\mathfrak{K}\) （§2 第 3 号命题 4）。对每个 \(\alpha \in \mathrm{A}\) ，映射 \(\Lambda\) 是标量本质 \(\mu_{\alpha}\) -可积的，我们令 \(\nu_{\alpha} = \int \lambda_t d\mu_{\alpha}(t)\) ；族 \((\nu_{\alpha})\) 可和，且其和等于 \(\nu\) （命题 1 的推论）。若 \(f\) 是定义在 X 上的正下半连续函数，则把证明的第一部分应用于测度 \(\mu_{\alpha}\) 与闭集 \(S_{\alpha}\) 就表明：

\(1^{\circ} h_{f}\) 是 \(\mu_{\alpha}\) -可测的（对每个 \(\alpha \in A\) ），从而是 \(\mu\) -可测的（§2 第 2 号命题 2），且

\[
2 ^ {\circ} \int^ {\bullet} f (x) d \nu_ {\alpha} (x) = \int^ {\bullet} d \mu_ {\alpha} (t) \int^ {\bullet} f (x) d \lambda_ {t} (x).
\]

对 \(\alpha\) 求和即得公式 (4)（ \(\S 2\) 第 2 号命题 1）。把前面的论证应用于任意测度 \(\mu' \leqslant \mu\) （这是合法的，因为 \(\Lambda\) 是标量本质 \(\mu'\) -可积且淡 \(\mu'\) -可测的，参看 \(\S 2\) 第 2 号命题 2），我们断定 \(\Lambda\) 是 \(\mu'\) -预适宜的，b) 得证。

最后，设 X 的拓扑有可数基，我们来证明每个标量本质 \(\mu\) -可积映射 \(\Lambda: t \mapsto \lambda_{t}\) （T 到 \(\mathcal{M}_{+}(X)\) 中的）都是淡 \(\mu\) -可测的。这将由下述引理得出：

引理. --- 设 X 是一个有可数基的局部紧空间。则在 \(\mathcal{K}(\mathrm{X})\) 中存在一个可数子集 S，它具有下列性质：对每个函数 \(f \in \mathcal{K}(\mathrm{X})\) ，存在 S 的元素的一个序列 \((f_n)\) 及一个正函数 \(\varphi \in S\) ，使对每个数 \(\varepsilon > 0\) ，只要 n 充分大就有 \(|f_n - f| \leqslant \varepsilon \varphi\) 。

设 \(X'\) 是 \(X\) 的 Alexandroff 紧化，它是一个可度量化紧空间（GT, IX, §2 第 9 号命题 16 及其推论）；我们把 \(\mathcal{K}(X)\) 等同于 \(\mathcal{C}(X')\) 的一个子集。设 \(S'\) 是 Banach 空间 \(\mathcal{C}(X')\) 的一个可数稠密子集（GT, X, §3 第 3 号定理 1）；可设 \(S'\) 包含常函数 \(n\) （对每个 \(n \in N\) ）。设 \((U_n)\) 是 \(X\) 中相对紧开集的一个序列，其并为 \(X\) ，使 \(\overline{U}_n \subset U_{n+1}\) 对一切 \(n\) 成立（GT, I, §9 第 9 号命题 15），并设 \(\varphi_n\) 是 \(\mathcal{K}_+(X)\) 中在 \(\overline{U}_n\) 上等于 1 的函数。我们用 \(S\) 表示 \(\mathcal{K}(X)\) 中形如 \(\varphi_n g\) （ \(n \in N, g \in S'\) ）的元素组成的可数集。若 \(f \in \mathcal{K}(X)\) ，设 \((g_n)\) 是 \(S'\) 的元素的一致收敛于 \(f\) 的一个序列，并设 \(k\) 是使 \(f\) 的支集含于 \(U_k\) 的一个整数。最后，设 \(m\) 是函数 \(g_n\) 的诸范数的一个上界整数。函数 \(f_n = \varphi_k g_n\) 属于 \(S\) 且满足引理的结论，其中 \(\varphi = m\varphi_k\) 。

引理既已建立，且映射 \(t \mapsto \lambda_t(g)\) 是本质 \(\mu\) -可积的（对每个 \(g \in S\) ），故映射 \(t \mapsto (\lambda_t(g))_{g \in S}\) （\(T\) 到 \(\mathbf{R}^S\) 中的）是 \(\mu\) -可测的（第四章，§5 第 3 号定理 1）。于是，集 \(\mathfrak{K}\) （由紧子集 \(K\) （\(T\) 中、使此映射在 \(K\) 上的限制连续者）组成）是 \(\mu\) -稠密的，且只需证明 \(\Lambda\) 在每个 \(K \in \mathfrak{K}\) 上的限制连续。设 \(f\) 是 \(\mathcal{H}(X)\) 的任意元素，\(f_n\) 与 \(\varphi\) 是 \(S\) 中满足引理结论的元素；则函数 \(t \mapsto \lambda_t(f)\) 在 \(K\) 上是连续函数 \(t \mapsto \lambda_t(f_n)\) 的一致极限；因而它在 \(K\) 上连续，命题得证。
% semantic-fragments: legacy-input-seam
\relax{}
